% ============================================================================
% Fichier  : partie2/P02_C03_sources_internationales.tex
% Titre    : Les Sources Internationales et Africaines
% Auteur   : MINKA MI NGUIDJOÏ Thierry Emmanuel
% Version  : 3.0 -- novembre 2026
% Corrigés : corriges/Corr_C03.tex
% ============================================================================
% Sources : lignes directrices de l'OCDE (1980, révisées en 2013) ; principes
% directeurs de l'ONU (résolution 45/95) ; Convention 108 et protocole STCE
% n° 223 ; Convention de Malabo ; cadre stratégique de l'UA en matière de
% données (2022) ; acte additionnel CEDEAO A/SA.1/01/10 ; Convention de
% Budapest. Tous paraphrasés. Statuts vérifiés en ligne le 12/11/2026 pour
% Malabo (en vigueur le 08/06/2023) et la Convention 108+ (34 ratifications
% sur 38), la directive CEMAC n° 07/08-UEAC-133-CM-18 et la loi ivoirienne
% n° 2013-450. Loi marocaine n° 09-08 et liste des Parties africaines à la
% Convention 108 : citées de mémoire, à confirmer.
% ============================================================================

\chapter{Les Sources Internationales et Africaines}
\label{chap:international}
\chaptermark{Les sources internationales}

\epigraph{%
    « Une loi nationale de protection des données\\
    est rarement une invention. C'est une traduction. »%
}{Principe du cours}

\niveauChapitre{Fondamental}{%
    Chapitres~\ref{chap:donnees} et~\ref{chap:principes}~; droit du numérique
    (UE MSN-2S3-JUR01)
}

\begin{boxobjectifs}
\begin{itemize}
    \item Retracer la filiation des principes de protection des données, des
          lignes directrices de l'OCDE à la loi n°~2024/017.
    \item Situer le droit à la vie privée dans les textes de droits
          fondamentaux qui lient le Cameroun.
    \item Donner le statut exact, pour le Cameroun, de chaque instrument
          africain et international.
    \item Dresser la fiche du droit d'un pays où vont des données.
    \item Classer un instrument comme opposable, contractuel ou comparatif
          avant de s'en servir dans un audit.
\end{itemize}
\end{boxobjectifs}

\bigskip

% ============================================================================
\section*{Avant de commencer : quatre pays, quatre droits}
% ============================================================================

Yaoundé, vendredi 3~juillet~2026. Dans la salle de réunion de KORO AUDIT
\& CONFORMITÉ, M.~NGANDO punaise une carte du monde au mur. Il y plante
quatre épingles pour les trois missions du cabinet. Abidjan, où l'éditeur
du logiciel bancaire de NYANGA CRÉDIT héberge les comptes. Casablanca, où
un centre d'appels répond la nuit aux clients de NOUN PAY. Francfort, où se
trouvent les serveurs de la même NOUN PAY. La Virginie, d'où lui parviennent
sa messagerie et sa mesure d'audience.

Un stagiaire demande quel droit s'applique à tout cela. « Le nôtre d'abord,
répond M.~NGANDO, puisque les personnes sont ici. Mais chaque épingle a aussi
son droit et certains de ces droits suivent les données jusqu'à nous. » Il
ajoute une cinquième épingle sur Strasbourg et une sixième sur Addis-Abeba.
« Celles-là n'hébergent rien. Elles expliquent pourquoi nos lois se
ressemblent. »

\begin{boxalert}
\textbf{Violation typique.} Citer dans un rapport d'audit un instrument
international comme s'il obligeait l'organisation, sans avoir vérifié que le
Cameroun l'a ratifié ni qu'il vise ce type d'acteur.

\textbf{Règle~:} avant de s'appuyer sur un texte, l'auditeur établit son
statut. Il est opposable, parce que le Cameroun y est Partie ou qu'une loi
le reprend. Il est contractuel, parce qu'un partenaire l'impose par contrat.
Ou il est comparatif et ne sert alors qu'à éclairer et à proposer.
\end{boxalert}

% ============================================================================
\section{La Filiation des Principes}
\label{sec:c03-filiation}
% ============================================================================

Les principes que le Chapitre~\ref{chap:principes} a lus dans les
articles~6 à~18 de la loi ont une histoire. Ils ont été formulés pour la
première fois à l'échelle internationale en 1980 et n'ont guère changé de
substance depuis. Les lois nationales les internalisent, chacune avec ses
ajouts.

\subsection{Les lignes directrices de l'OCDE}

En 1980, l'Organisation de coopération et de développement économiques
adopte des lignes directrices sur la protection de la vie privée et les flux
transfrontières de données~\cite{ocde1980}. Elles ne sont pas un traité. Ce
sont des recommandations, révisées en 2013, que les États membres
s'engagent à suivre. Elles énoncent huit principes qui forment encore le
socle de presque toutes les lois du monde.

\begin{tableaucours}{L{4.2cm}L{4.6cm}L{3.2cm}}{%
    \tableauHeader{Principe de l'OCDE (paraphrasé)} &
    \tableauHeader{Idée} &
    \tableauHeader{Loi n°~2024/017}
}
Limitation de la collecte & Collecter loyalement, avec le consentement ou
l'information de la personne & Art.~9, 10 \\
Qualité des données & Des données exactes, complètes et à jour & Art.~12 \\
Spécification des finalités & Dire pourquoi l'on collecte, au plus tard à la
collecte & Art.~11, 21 \\
Limitation de l'utilisation & Ne pas détourner les données de leur finalité
& Art.~11, 67 \\
Garanties de sécurité & Protéger contre la perte, l'accès et la divulgation
& Art.~7, 15, 27 \\
Transparence & Faire connaître l'existence et la nature des traitements &
Art.~14, 29 \\
Participation individuelle & Permettre à la personne d'accéder à ses données
et de les contester & Art.~37 à~46 \\
Responsabilité & Rendre le responsable comptable du respect des principes &
Art.~20, 54 à~71 \\
\end{tableaucours}
\vspace{-0.8em}
\begin{center}{\small\itshape Tableau~3.1. Les huit principes de l'OCDE et
leur écho dans la loi}\end{center}

Le Cameroun n'est pas membre de l'OCDE. Le tableau ne décrit donc pas une
obligation, il décrit une parenté. Un auditeur qui la connaît lit plus vite
une loi étrangère et repère plus vite ce que la sienne a de particulier.

\subsection{Les principes directeurs des Nations Unies}

L'Assemblée générale des Nations Unies a adopté en 1990 des principes
directeurs pour la réglementation des fichiers informatisés contenant des
données à caractère personnel~\cite{onu1990}. Le texte reprend la licéité, la
loyauté, l'exactitude, la finalité, l'accès et la sécurité. Il ajoute deux
idées que l'on retrouve dans la loi camerounaise~: la non-discrimination, à
l'article~17 et l'existence d'une autorité de contrôle indépendante, à
l'article~53. C'est une résolution, qui recommande sans lier.

\subsection{La Convention 108 et sa version modernisée}

La Convention du Conseil de l'Europe du 28~janvier~1981, dite
Convention~108~\cite{coe1981conv108}, est le premier traité contraignant
consacré à la protection des données. Elle est ouverte aux États qui ne sont
pas membres du Conseil de l'Europe et plusieurs États africains y ont
adhéré, parmi lesquels le Maroc, le Sénégal, la Tunisie, Maurice et le
Cap-Vert. Le Cameroun n'y est pas Partie.

Un protocole de 2018 la modernise~: c'est la Convention~108+~\cite{coe2018conv108plus}.
Il ajoute la notification des violations, la protection des données dès la
conception et des droits face aux décisions automatisées. Il n'entrera en
vigueur qu'à la trente-huitième ratification. En 2026, le Conseil de
l'Europe en comptait trente-quatre.

\begin{boxinfo}
\textbf{Pourquoi un traité européen intéresse l'auditeur camerounais.} Le
centre d'appels de NOUN PAY est au Maroc, qui est Partie à la
Convention~108. Cette adhésion est un fait vérifiable sur le site du
dépositaire. Elle ne vaut pas reconnaissance par l'Autorité camerounaise,
mais elle entre dans la description du pays de destination que l'auditeur
joint au dossier de transfert (section~\ref{sec:c03-pays}).
\end{boxinfo}

% ============================================================================
\section{Le Droit à la Vie Privée}
\label{sec:c03-vieprivee}
% ============================================================================

La protection des données procède d'un droit plus ancien, le droit au
respect de la vie privée. Quatre textes le portent pour le Cameroun.

\begin{tableaucours}{L{4.4cm}L{4.4cm}L{3.2cm}}{%
    \tableauHeader{Texte} &
    \tableauHeader{Ce qu'il protège} &
    \tableauHeader{Statut}
}
Déclaration universelle des droits de l'homme (1948), article~12 & Vie
privée, famille, domicile, correspondance & Référence, visée par le préambule
de la Constitution \\
Pacte international relatif aux droits civils et politiques (1966),
article~17 & Les mêmes, contre les immixtions arbitraires ou illégales &
Traité, opposable \\
Constitution, préambule & Inviolabilité du domicile, secret de toute
correspondance & Norme suprême \\
Charte africaine des droits et du bien-être de l'enfant (1990), article~10 &
Vie privée de l'enfant & Traité, opposable \\
\end{tableaucours}
\vspace{-0.8em}
\begin{center}{\small\itshape Tableau~3.2. Le droit à la vie privée dans les
textes de droits fondamentaux}\end{center}

La Charte africaine des droits de l'homme et des peuples de 1981 ne contient
pas d'article consacré à la vie privée. Le fait est souvent ignoré. En
Afrique, la protection des données s'est donc construite par des textes
spécialisés plutôt que par la Charte.

Ces textes ne fournissent pas de critères d'audit~: ils sont trop généraux
pour être cotés. Ils servent au juge, qui peut y rattacher une atteinte que
la loi spéciale ne prévoit pas et ils rappellent à l'auditeur que la
personne concernée est titulaire d'un droit fondamental avant d'être une
ligne dans un registre.

% ============================================================================
\section{Le Cadre Africain}
\label{sec:c03-afrique}
% ============================================================================

\subsection{La Convention de Malabo}

La Convention de l'Union africaine sur la cybersécurité et la protection des
données à caractère personnel a été adoptée à Malabo le
27~juin~2014~\cite{ua2014malabo}. Elle traite trois sujets~: les transactions
électroniques, la protection des données et la cybersécurité. Son chapitre~II
demande à chaque État de se doter d'un cadre légal, d'une autorité nationale
indépendante, de formalités préalables aux traitements, de principes et de
droits proches de ceux du Tableau~3.1.

Elle est entrée en vigueur le 8~juin~2023, trente jours après le dépôt du
quinzième instrument de ratification. Le Cameroun l'a signée et ne l'avait
pas ratifiée à la date de ce cours. Elle ne l'oblige donc pas encore.

Sa lecture éclaire pourtant deux choix de la loi n°~2024/017. Le régime
d'autorisation préalable de l'article~19 prolonge les formalités préalables
de la Convention. La liste ouverte des données sensibles reprend l'esprit de
la sienne. Le législateur camerounais a internalisé un texte qu'il n'avait
pas ratifié.

\subsection{Le cadre stratégique de l'Union africaine}

En 2022, l'Union africaine a adopté un cadre stratégique en matière de
données~\cite{ua2022donnees}. Ce n'est pas un texte normatif. Il fixe une
orientation~: des données qui circulent entre États africains, des autorités
qui coopèrent, des règles nationales qui convergent. L'auditeur le cite
quand une organisation invoque la souveraineté des données pour refuser tout
hébergement hors du pays, ou au contraire le marché africain pour justifier
un transfert. Dans les deux cas, le cadre donne le vocabulaire et non la
règle.

\subsection{Les textes régionaux}

L'Afrique de l'Ouest a pris de l'avance. L'acte additionnel de la CEDEAO du
16~février~2010~\cite{cedeao2010acte} oblige ses États membres à légiférer
et la plupart l'ont fait. Le Cameroun n'appartient pas à la CEDEAO~; l'acte
ne lui est pas applicable.

En Afrique centrale, il n'existe pas de texte équivalent consacré aux
données personnelles. La CEMAC a adopté le 19~décembre~2008 cinq directives sur les
communications électroniques. L'une d'elles, la directive
n°~07/08-UEAC-133-CM-18, protège les droits des utilisateurs et demande aux
États de garantir la confidentialité des communications et des données de
trafic~\cite{cemac2008}. La loi camerounaise n°~2010/013 l'a transposée. Le secteur bancaire et celui des paiements
relèvent en outre de règlements communautaires et de la Commission bancaire
de l'Afrique centrale, qui encadrent le secret et la conservation des
données des clients (Chapitre~\ref{chap:loi2010}). Pour une microfinance ou
une fintech, ces textes sont opposables et s'appliquent directement dans
leur domaine.

\begin{boxinfo}
\textbf{Les autorités se parlent.} Les autorités africaines de protection
des données se sont réunies en 2016 dans un réseau, le RAPDP. Quand
l'Autorité camerounaise sera organisée, c'est là qu'elle trouvera ses
homologues. La loi lui donne mission de coopérer avec elles (article~53).
\end{boxinfo}

% ============================================================================
\section{Le Droit des Pays où Vont les Données}
\label{sec:c03-pays}
% ============================================================================

L'article~32 de la loi subordonne tout transfert à la vérification que le
pays de destination « offre un niveau de protection suffisant ». C'est à
l'Autorité de l'apprécier et le Chapitre~\ref{chap:transferts} dira comment
coter tant qu'elle ne l'a pas fait. L'organisation, elle, doit préparer la
description du pays. L'auditeur vérifie que cette description existe et
qu'elle repose sur des sources primaires.

\begin{boxprocedure}
\textbf{Dresser la fiche d'un pays de destination.}
\begin{enumerate}
    \item La loi générale de protection des données~: référence, date, champ.
    \item L'autorité de contrôle~: nom, indépendance, pouvoirs de sanction.
    \item Les engagements internationaux~: Convention~108, Convention de
          Malabo, autres.
    \item Les reconnaissances reçues d'autres autorités, par exemple une
          décision d'adéquation de l'Union européenne.
    \item L'accès des autorités publiques aux données et les recours ouverts
          aux étrangers.
    \item Les sources consultées et la date de consultation.
\end{enumerate}
\end{boxprocedure}

\begin{tableaucours}{L{2.2cm}L{3.6cm}L{2.8cm}L{3.4cm}}{%
    \tableauHeader{Pays} &
    \tableauHeader{Loi générale} &
    \tableauHeader{Autorité} &
    \tableauHeader{Engagements}
}
Côte d'Ivoire & Loi n°~2013-450 du 19~juin~2013~\cite{ci2013loi} & Autorité de régulation des télécommunications
(ARTCI) & Malabo ratifiée~; CEDEAO \\
Maroc & Loi n°~09-08 de 2009~\cite{ma2009loi} & Commission nationale
(CNDP) & Convention~108 \\
Union européenne & RGPD (Chapitre~\ref{chap:rgpd}) & Une autorité par État,
un comité européen & Convention~108 \\
États-Unis & Pas de loi fédérale générale~; lois sectorielles et lois
d'États & Pas d'autorité unique & Aucun traité de protection des données \\
\end{tableaucours}
\vspace{-0.8em}
\begin{center}{\small\itshape Tableau~3.3. Les quatre destinations des
données des trois missions (à vérifier sur les textes avant tout usage)}\end{center}

La dernière ligne mérite un mot. Les États-Unis n'ont pas de loi générale
comparable aux trois autres. La protection y dépend du secteur (santé,
finance, enfants) et de l'État (la Californie a la loi la plus connue). Les
entreprises américaines s'appuient souvent sur des cadres volontaires, dont
celui du NIST, que le Chapitre~\ref{chap:normes} présente. Une fiche pays
honnête l'écrit sans juger~: l'absence de loi générale est un fait, pas une
cote.

Aucun pays d'Afrique subsaharienne ne bénéficiait en 2026 d'une décision
d'adéquation de l'Union européenne. Les données qui partent d'Europe vers le
Cameroun voyagent donc sous clauses contractuelles, ce qui aura des
conséquences pour les sous-traitants camerounais de clients européens
(Chapitre~\ref{chap:rgpd}).

% ============================================================================
\section{La Coopération Pénale}
\label{sec:c03-budapest}
% ============================================================================

Les données franchissent aussi les frontières à la demande des enquêteurs.
La Convention de Budapest sur la cybercriminalité de
2001~\cite{budapest2001} organise cette coopération~: conservation rapide de
données stockées à l'étranger, entraide, points de contact disponibles en
permanence. Le Cameroun y est Partie depuis le 15~décembre~2023 et elle est
entrée en vigueur à son égard le 1\textsuperscript{er}~avril~2024, d'après
les travaux de l'auteur et de ses coauteurs~\cite{minka2026chainepenale}.

Pour l'auditeur de la protection des données, la Convention a deux effets.
Elle fonde des demandes que l'organisation auditée peut recevoir et qu'elle
doit savoir traiter sans violer l'article~32. Elle rencontre aussi le régime
d'autorisation préalable des transferts, rencontre que le
Chapitre~\ref{chap:recherche} étudie comme une tension transfrontalière.

% ============================================================================
\section{Le Statut de Chaque Instrument}
\label{sec:c03-statut}
% ============================================================================

\begin{tableaucours}{L{3.6cm}L{4.2cm}L{4.2cm}}{%
    \tableauHeader{Instrument} &
    \tableauHeader{Statut pour le Cameroun} &
    \tableauHeader{Usage dans l'audit}
}
Pacte de 1966, article~17 & Opposable & Fondement général, non coté \\
Convention de Budapest (2001) & Opposable depuis 2024 & Demandes des
autorités, coopération pénale \\
Textes CEMAC et COBAC & Opposables dans leur domaine & Critères sectoriels
(banque, paiement, télécommunications) \\
Convention de Malabo (2014) & Signée, non ratifiée & Comparatif et
prospectif \\
Convention~108 et 108+ & Non Partie & Comparatif~; fiche des pays Parties \\
Acte additionnel CEDEAO (2010) & Non applicable & Comparatif~; fiche des
pays de la CEDEAO \\
Lignes directrices de l'OCDE & Non membre & Comparatif \\
Principes directeurs de l'ONU (1990) & Recommandation & Comparatif \\
Cadre de l'UA sur les données (2022) & Orientation politique & Vocabulaire \\
RGPD (UE, 2016) & Voir Chapitre~\ref{chap:rgpd} & Opposable à certains
acteurs, contractuel pour d'autres \\
\end{tableaucours}
\vspace{-0.8em}
\begin{center}{\small\itshape Tableau~3.4. Les instruments internationaux et
leur statut}\end{center}

Trois statuts suffisent à l'auditeur.

\begin{itemize}
    \item \textbf{Opposable.} Le texte lie l'organisation. Un écart se cote
          \NON{} et se rattache à l'article.
    \item \textbf{Contractuel.} Le texte ne lie pas l'organisation par
          lui-même, mais un contrat le lui impose. L'écart est un manquement
          au contrat, que l'auditeur cote si le contrat figure parmi ses
          critères.
    \item \textbf{Comparatif.} Le texte n'oblige pas. Il sert à interpréter,
          à combler un silence par un critère de substitution déclaré, ou à
          proposer une bonne pratique. Il ne fonde jamais un \NON{} à lui
          seul.
\end{itemize}

Un statut change. La ratification de Malabo ferait passer la Convention de
comparative à opposable. L'auditeur date donc le tableau et le revoit à
chaque mission.

% ============================================================================
\section{La Carte des Trois Missions}
\label{sec:c03-cas}
% ============================================================================

\begin{boxscenario}{Missions NYANGA, MBOLO et NOUN PAY. La carte de M. NGANDO, 3 juillet 2026}
\textbf{Constatations.}
\begin{itemize}
    \item NYANGA CRÉDIT~: un flux vers la Côte d'Ivoire (hébergement du
          logiciel bancaire). Aucune description du droit ivoirien au
          dossier. Établissement soumis aux textes de la COBAC.
    \item CENTRE MÉDICAL MBOLO~: aucun flux hors du Cameroun identifié à ce
          stade. Prestataire à Douala.
    \item NOUN PAY~: flux vers l'Union européenne, les États-Unis et le
          Maroc, un quatrième vers un pays inconnu (les selfies). Aucune
          fiche pays. Établissement de paiement soumis aux textes
          communautaires.
\end{itemize}

\textbf{Analyse} (confiance élevée sur les flux déclarés, faible sur leur
exhaustivité). Aucune des trois organisations n'a décrit le droit des pays
où vont ses données. Deux d'entre elles relèvent de textes communautaires
qu'elles n'ont pas rapprochés de la loi n°~2024/017.
\end{boxscenario}

\begin{boxlecture}{Le dossier international de NOUN PAY}
\begin{tableaucours}{L{7.6cm}cL{3.4cm}}{%
    \tableauHeader{Proposition} &
    \tableauHeader{Cote} &
    \tableauHeader{Pièce}
}
Les pays de destination sont identifiés & \NON{} & Selfies, pays inconnu \\
Une fiche existe pour chaque pays de destination & \NON{} & Aucune fiche \\
Les fiches reposent sur des sources primaires datées & \NON{} & Sans objet,
aucune fiche \\
Les textes communautaires applicables sont recensés & \NON{} & Aucune
liste \\
Le Maroc est Partie à la Convention~108 & \OUI{} & Liste du dépositaire \\
Le niveau de protection de chaque pays est reconnu suffisant & \NEK{} &
Liste de l'Autorité non publiée \\
\end{tableaucours}

\medskip
\textbf{Décompte.} Un \OUI{}, quatre \NON{}, un \NEK{} sur six.
\hfill\lecture{0{,}17}{0{,}17}{1{,}25}

\medskip
\textbf{Lecture.} Le seul \OUI{} ne doit rien à l'entreprise~: c'est un fait
de droit international que l'auditeur a vérifié lui-même. Le \NEK{} est de
critère. Les quatre \NON{} relèvent tous de la préparation du dossier, qui
dépend de NOUN PAY seule.
\end{boxlecture}

% ============================================================================
\section{Ce que les Sources demandent à chaque Métier}
\label{sec:c03-metiers}
% ============================================================================

\begin{boxmetier}{Juriste conformité}
Tenir le tableau des instruments avec leur statut et leur date de
vérification. Aller au dépositaire~: l'Union africaine pour Malabo, le
Conseil de l'Europe pour la Convention~108 et pour Budapest.
\end{boxmetier}

\begin{boxmetier}{Auditeur de la protection des données}
Écrire dans le plan d'audit le statut de chaque texte retenu comme critère.
Ne jamais coter \NON{} sur la seule foi d'un texte comparatif.
\end{boxmetier}

\begin{boxmetier}{Délégué à la protection des données}
Tenir une fiche par pays de destination et la joindre au registre. La
mettre à jour quand un prestataire change de région d'hébergement.
\end{boxmetier}

\begin{boxresume}
\begin{itemize}
    \item Les principes de la loi n°~2024/017 descendent des huit principes
          de l'OCDE de 1980, repris par l'ONU, le Conseil de l'Europe et
          l'Union africaine.
    \item Le droit à la vie privée est garanti par le Pacte de 1966 et par
          le préambule de la Constitution~; la Charte africaine de 1981 ne
          le mentionne pas.
    \item La Convention de Malabo est en vigueur depuis le 8~juin~2023. Le
          Cameroun l'a signée sans la ratifier~: elle reste comparative.
    \item Un instrument est opposable, contractuel ou comparatif. Seuls les
          deux premiers fondent une cote \NON{}.
    \item L'organisation qui transfère des données tient une fiche du droit
          de chaque pays de destination, fondée sur des sources primaires.
\end{itemize}
\end{boxresume}

\begin{boxinfo}
\textbf{Points de contrôle.} Vous avez compris ce chapitre si vous savez~:
\begin{itemize}
    \item rattacher un article de la loi à un principe de l'OCDE~;
    \item dire pourquoi la Convention de Malabo ne fonde pas un constat
          d'audit au Cameroun~;
    \item dresser la fiche d'un pays de destination.
\end{itemize}
\end{boxinfo}

\begin{boxexo}[title={\ding{45}\quad Exercice~3.1. Quel statut \niveaubase}]
Donnez le statut (opposable, contractuel, comparatif) de chaque texte pour
NYANGA CRÉDIT~: (a) la Convention de Malabo~; (b) le Pacte de 1966~;
(c) l'acte additionnel de la CEDEAO~; (d) un règlement de la COBAC~;
(e) les lignes directrices de l'OCDE.
\end{boxexo}

\begin{boxexo}[title={\ding{45}\quad Exercice~3.2. La filiation \niveaumoyen}]
Pour chacun des articles suivants de la loi n°~2024/017, nommez le principe
de l'OCDE qu'il internalise, puis dites ce que la loi camerounaise y ajoute~:
(a) article~9~; (b) article~13~; (c) article~19~; (d) article~27.
\end{boxexo}

\begin{boxexo}[title={\ding{45}\quad Exercice~3.3. La fiche pays \niveaumoyen}]
Dressez, en suivant les six rubriques de la procédure, la fiche de la Côte
d'Ivoire pour le dossier de transfert de NYANGA CRÉDIT. Indiquez pour chaque
rubrique la source primaire que vous consulteriez et ce que vous écrivez si
vous ne la trouvez pas.
\end{boxexo}

\begin{boxexo}[title={\ding{45}\quad Exercice~3.4. Si le Cameroun ratifiait \niveauexpert}]
Supposez que le Cameroun dépose son instrument de ratification de la
Convention de Malabo.
\begin{enumerate}[(a)]
    \item À quelle date la Convention entrerait-elle en vigueur à son égard~?
          Quelle source consultez-vous pour le savoir~?
    \item Quelles lignes du Tableau~3.4 changent~?
    \item Un audit en cours doit-il modifier ses critères~? Rédigez la note
          au commanditaire.
\end{enumerate}
\end{boxexo}

\bigskip

\begin{boxinfo}
\textbf{Pour aller plus loin.}
\begin{itemize}
    \item Les lignes directrices de l'OCDE~\cite{ocde1980}, pour la
          formulation d'origine des huit principes.
    \item La Convention de Malabo~\cite{ua2014malabo}, chapitre~II et sa
          liste d'état tenue par l'Union africaine.
    \item La Convention~108+~\cite{coe2018conv108plus} et son rapport
          explicatif.
\end{itemize}
\end{boxinfo}

\begin{boxinfo}
\textbf{Transition.} Parmi tous ces textes, un seul suit les données
au-delà de ses frontières avec assez de force pour atteindre une entreprise
de Bafoussam~: le règlement européen. C'est l'objet du
Chapitre~\ref{chap:rgpd}.
\end{boxinfo}
